% ============================================================
% Capitolo 4 — Metodologia sperimentale (traduzione italiana di
% chapters/ch04_methodology.tex; label invariati)
% ============================================================
\chapter{Metodologia sperimentale}
\label{ch:methodology}

La piattaforma del \cref{ch:system} non è stata progettata in un solo
passaggio: è il prodotto di una campagna iterativa in cui ogni modifica
doveva superare un gate quantitativo predefinito prima di entrare nel
trunk. Questo capitolo documenta le regole di misura
(\cref{sec:metrics}), il protocollo di gate (\cref{sec:gate}), i difetti
di validità trovati e corretti lungo il percorso (\cref{sec:bugs}), una
sintesi delle iterazioni stesse (\cref{sec:iterations}) e il design
della campagna di confronto finale (\cref{sec:campaign}).

% ------------------------------------------------------------
\section{Metriche e regole di misura}
\label{sec:metrics}

\paragraph{Unità di misura.}
L'unità di analisi è il \emph{record a livello di agente}: una voce JSON
per agente per episodio (sei per episodio, \cref{sec:termination}).
Tutti gli aggregati sono calcolati sui record dopo deduplica per
identificatore di episodio e di agente (in valutazione, lo scenario fa
parte della chiave, poiché gli identificatori di episodio possono
ripetersi tra scenari). Due condizioni di integrità sono imposte prima
che qualunque numero venga riportato: esattamente sei record per
episodio, e nessun valore non finito in alcun campo. Una run che fallisce
una delle due condizioni non viene analizzata oltre.

\paragraph{Criterio di successo rigoroso.}
Un agente ha successo se e solo se la sua ragione di terminazione è
\texttt{route\_complete}. In particolare, \texttt{route\_short} ---
raggiungere la fine di un percorso degenere troppo corto --- \emph{non}
è successo (\cref{sec:termination}). Questa regola è stata messa alla
prova durante lo sviluppo: una modifica che contava \texttt{route\_short}
come successo è stata introdotta e successivamente revertita
(\cref{sec:iterations}), mantenendo il criterio rigoroso al costo di
numeri assoluti più bassi. Il \cref{ch:discussion} discute perché tassi
di successo assoluti di fascia media siano un pregio, non un difetto,
della misura risultante: lontano dal pavimento e dal soffitto, la
metrica conserva spazio per muoversi in entrambe le direzioni, che è
esattamente ciò di cui uno studio comparativo ha bisogno.

\paragraph{Tre viste di reporting.}
Ogni risultato è riportato per veicoli e pedoni separatamente, più una
vista combinata; le due classi hanno task diversi, spazi di azione
diversi e limiti strutturali diversi (\cref{sec:ped-limits}), e il
registro sperimentale mostra effetti di regime che esistono per una sola
classe. Il successo congiunto a livello di episodio (``tutti e sei gli
agenti riescono'') non è mai usato: con agenti correlati in un mondo
condiviso compone i fallimenti in modo moltiplicativo e nasconde la
struttura specifica di classe.

\paragraph{Catalogo delle metriche.}
La \cref{tab:metrics} definisce le metriche usate in tutto il
\cref{ch:results}. I valori cumulativi su tutti gli episodi di
addestramento sono l'aggregazione primaria; il successo a finestra
scorrevole (finestra 50) è usato solo dal gate del curriculum
(\cref{sec:regimes}); la dinamica entro la run è mostrata dividendo
l'addestramento in quartili (Q1--Q4).

\begin{table}[t]
  \centering\small
  \caption{Catalogo delle metriche. Tutti i tassi sono quote dei record
  a livello di agente.}
  \label{tab:metrics}
  \begin{tabular}{@{}p{3.6cm}p{9.6cm}@{}}
    \toprule
    Metrica & Definizione \\
    \midrule
    Tasso di successo (SR) & record con \texttt{route\_complete} \\
    Tasso di collisione & record con \texttt{collision} \\
    Tasso di offroad & record con \texttt{offroad} (veicoli) \\
    Tasso di stuck & record troncati con completamento $< 0.3$ o rilevatore di loop attivo \\
    Tasso di timeout & record troncati con completamento $\geq 0.3$ \\
    Stuck$+$timeout & somma dei due precedenti; target primario anti-stallo \\
    Completamento del percorso & frazione di waypoint del percorso consumati, $[0,1]$ \\
    Efficienza di percorso & lunghezza ottima completata / distanza effettivamente percorsa, con tetto a 1; definita 0 per i record stuck \\
    Velocità media & km/h per i veicoli, m/s per i pedoni \\
    Passi di non-progresso & passi dall'ultimo avanzamento di waypoint, a fine episodio \\
    Diagnostica dei percorsi & quote di \texttt{route\_source}, flag under-target e too-short, conteggi di attraversamenti (\cref{sec:routes}) \\
    \bottomrule
  \end{tabular}
\end{table}

% ------------------------------------------------------------
\section{Protocollo di sviluppo guidato da gate}
\label{sec:gate}

La varianza da run a run in questo ambiente è grande (quantificata più
avanti in $\sim$10 punti percentuali di SR veicoli tra repliche
identiche), e il budget di calcolo disponibile non permetteva di
replicare ogni modifica candidata. Il protocollo adottato in risposta ha
quattro regole.

\paragraph{1. Una manopola alla volta.}
Ogni candidato è un diff stretto e isolato comportamentalmente --- un
singolo termine di reward, una singola soglia, un singolo gruppo di
osservazioni. I candidati che toccano lo stesso file sono valutati in
sequenza, mai impilati. Quando una modifica in bundle ha fallito il suo
gate, l'attribuzione tra le sue metà è stata trattata come irrisolta
anziché indovinata.

\paragraph{2. A/B appaiato a budget diagnostico fisso.}
Ogni candidato gira per $3\times10^{5}$ passi contro la baseline
corrente del trunk, con lo stesso seed di esperimento e, grazie al
seeding deterministico dei percorsi della \cref{sec:seeding}, una
sequenza di percorsi appaiata alla baseline per costruzione. Le run
della fase di sviluppo erano bloccate sul livello easy salvo quando il
candidato riguardava il curriculum stesso, così che il mescolamento dei
livelli non potesse confondere il confronto. Quando si confrontano run
di lunghezza diversa, i record sono troncati a conteggi di episodi
uguali.

\paragraph{3. Un gate quantitativo predefinito.}
Per le modifiche orientate ai veicoli il candidato è promosso solo se,
rispetto alla baseline appaiata: l'SR veicoli migliora di almeno
$+2.0$\,pp; stuck$+$timeout migliora di almeno $-2.0$\,pp; collisione e
offroad peggiorano ciascuno di non più di $+1.0$\,pp; ed entrambe le
condizioni di integrità della \cref{sec:metrics} valgono. Le modifiche
che toccano macchinari condivisi portano una condizione laterale
\emph{must-not-collapse} sull'altra classe (SR pedoni entro $-2$\,pp).
Un gate fallito significa che il candidato viene revertito; il trunk non
è mai lasciato in uno stato parzialmente valutato.

\paragraph{4. Meccanismo ed esito sono registrati separatamente.}
Per ogni candidato il registro annota se il \emph{meccanismo} inteso si
è verificato (per es.\ la varianza spiegata del critic che sale da
$\approx 0$ a $0.87$ dopo una riparazione della value loss) e se il
\emph{gate comportamentale} è passato --- i due spesso dissentono. Un
piccolo numero di manopole che hanno fallito il gate è stato
consapevolmente mantenuto con l'ipotesi registrata come falsificata
(per es.\ una penalità di collisione decuplicata che ha lasciato piatto
il tasso di collisione), perché revertirle avrebbe ripristinato un
difetto o perché il fallimento stesso era informativo. Queste eccezioni
sono esplicite nel registro (\cref{app:registry}); nulla è stato
mantenuto silenziosamente.

Le traiettorie per quartile entro la run completano il gate: un
candidato i cui numeri cumulativi passano ma il cui Q4 declina (una
firma osservata per un retuning del passo dei pedoni) è trattato come
rigettato.

% ------------------------------------------------------------
\section{Validità dell'ambiente: bug trovati e corretti}
\label{sec:bugs}

In una piattaforma multi-agente custom l'apparato di misura è parte
dell'esperimento. Tre episodi durante lo sviluppo hanno prodotto
apparenti effetti di policy che erano in realtà artefatti; sono
riportati qui sia perché hanno plasmato il protocollo finale sia perché
le loro firme di rilevamento possono essere utili ad altri.

\paragraph{Bug della lunghezza dei percorsi (distribuzione dei task).}
Il route planner poteva emettere percorsi molto più lunghi del target di
livello, quindi episodi nominalmente ``easy'' spesso non erano facili.
Imporre i limiti $[0.5, 2.0]\times$ il target (\cref{sec:routes}) ha
spostato l'SR cumulativo di addestramento dei veicoli di $\approx
+50$\,pp \emph{senza alcuna modifica alla policy o alla reward} --- un
puro cambiamento di distribuzione dei task. Ne sono state adottate due
conseguenze: le metriche assolute sono confrontabili solo entro la
stessa versione dell'ambiente, e tutti i risultati precedenti sono stati
rietichettati come specifici della loro era (il pavimento dei veicoli
del pilota di aprile, \cref{sec:pilot}, è stato poi ricondotto a questo
difetto). Questo episodio è l'evidenza interna più forte per una regola
usata in tutta questa tesi: \emph{mai confrontare numeri assoluti tra
versioni dell'ambiente; ancorare le affermazioni a contrasti appaiati
entro una versione}.

\paragraph{Bug del seeding dei percorsi (appaiamento rotto).}
I seed dei percorsi per agente derivavano originariamente dalla
\texttt{hash()} built-in di Python, che è salata per processo: due run
nominalmente identiche estraevano sequenze di percorsi \emph{diverse},
degradando silenziosamente ogni confronto A/B. La correzione deriva i
seed da una \texttt{SeedSequence} esplicita su seed del traffico,
contatore di reset e un digest stabile dell'identificatore dell'agente
(\cref{sec:seeding}). La lezione: l'appaiamento va costruito e
verificato (statistiche di generazione dei percorsi byte-identiche tra i
bracci), non assunto.

\paragraph{Bug della pipeline di valutazione (collasso della misura).}
La prima valutazione di checkpoint ha restituito un tasso di successo
esattamente del 33.33\% in ogni scenario --- uno dei tre veicoli a
completare, ogni episodio. Questo pattern ``troppo pulito per essere
vero'' ha esposto quattro difetti accoppiati: (i) il seed di scenario
per episodio era costante, quindi ogni episodio rigiocava la stessa
configurazione di traffico; (ii) la policy era valutata sulla sua media
deterministica, collassando la policy gaussiana su una singola
traiettoria per scenario; (iii) il sottoprocesso di valutazione non
inizializzava mai il logger degli episodi, quindi nessun record a
livello di agente veniva scritto; e (iv) il seed torch era costante tra
i sottoprocessi degli episodi, quindi anche azioni stocastiche avrebbero
rigiocato flussi di rumore correlati. Tutti e quattro sono stati
corretti (seed di scenario e campionatore per episodio, valutazione
stocastica, propagazione del logger con riavvio pulito), la valutazione
deterministica parziale è stata scartata e il protocollo stocastico
della \cref{sec:eval-pipeline} è diventato lo standard. Nessuna di
queste correzioni tocca l'addestramento; per le regole stesse del
progetto, nessuna affermazione di policy è derivata da esse.

La lezione generale di questa sezione: la validità della misura va
stabilita \emph{prima} del tuning, e risultati implausibilmente regolari
sono un segnale di rilevamento, non una buona notizia.

% ------------------------------------------------------------
\section{Sintesi dello sviluppo iterativo}
\label{sec:iterations}

La \cref{tab:candidates} condensa il registro di sviluppo; il registro
completo per candidato, con run ID e meccanismi, è
nell'\cref{app:registry}. Lo sviluppo è proceduto in sei fasi:

\begin{enumerate}
  \item \textbf{Pilota architetturale (aprile 2026).} Nove run che
        incrociano encoder del critic e regimi di addestramento,
        nell'era dell'ambiente pre-correzioni. I dati erano integri ma i
        veicoli sedevano a un SR di pavimento (0.2--3.8\%) in ogni
        cella, rendendo il confronto architetturale non informativo ---
        e reindirizzando lo sforzo verso la validità dell'ambiente
        (\cref{sec:pilot}).
  \item \textbf{Correttezza.} Le correzioni di lunghezza e seeding dei
        percorsi della \cref{sec:bugs}, più strumentazione diagnostica
        per passo (progresso, passi di non-progresso, cause di stuck,
        provenienza dei percorsi).
  \item \textbf{Osservabilità.} Feature veicolo consapevoli del
        percorso, poi l'estensione a 47D che espone lo stato di
        non-progresso, lo stato della penalità anti-loop e il tempo
        rimanente (\cref{sec:obs-actions}), chiudendo il gap di
        aliasing di stato tra reward e osservazione. La baseline
        easy-locked risultante era sana (plateau SR $\approx 79$\%,
        varianza spiegata del critic 0.983).
  \item \textbf{Iterazioni su ottimizzatore e reward.} Una sequenza di
        candidati a manopola singola con verdetti misti
        (\cref{tab:candidates}); in particolare, una riparazione della
        funzione di valore il cui meccanismo è stato confermato ma il
        cui gate è fallito, e una penalità di collisione decuplicata che
        ha lasciato piatte le collisioni --- l'evitamento delle
        collisioni non era apprendibile dalle osservazioni a 44D di
        quell'era, un limite di percezione più che un problema di peso
        della reward.
  \item \textbf{Accoppiamento tra classi.} Pedoni più veloci saturano il
        canale di hazard dei veicoli e possono spingere i veicoli in un
        equilibrio difensivo a velocità zero attraverso il gate
        \texttt{safe\_to\_push}. La risoluzione promossa (V1) ha alzato
        la soglia del gate da 0.75 a 0.85, mantenendo attiva l'urgenza
        dei veicoli sotto hazard moderato; un retuning del passo dei
        pedoni nella direzione opposta (V3) è stato rigettato per una
        traiettoria dell'ultimo quartile in declino.
  \item \textbf{Meccanica del curriculum e trunk finale.} Sblocco a
        finestra con soglie di ripasso, il gate sulla policy più debole
        e cap di pressione rialzati (\cref{sec:regimes}); infine, il
        tetto agli attraversamenti pedonali (\texttt{max\_cross} $8 \to
        3$), l'ultima modifica entrata nel trunk prima della campagna
        (SR veicoli $+9.2$\,pp, collisione $-2.3$\,pp, offroad
        $-2.2$\,pp, pedoni piatti).
\end{enumerate}

Tre repliche a lunghezza piena ($3\times10^{6}$ passi) dell'identico
trunk pre-campagna quantificano la varianza da run a run: l'SR
cumulativo dei veicoli ha spaziato tra 52.2 e 62.4\% sotto lo stesso
seed --- CARLA non è deterministico sotto carico nemmeno se seedato.
Questa banda di $\sim$10\,pp è il metro contro cui ogni effetto del
\cref{ch:results} viene letto, e motiva sia il design appaiato sia la
cautela attribuita ai delta a replica singola (\cref{sec:threats}).

\begin{table}[t]
  \centering\small
  \caption{Registro condensato dei candidati (fase di sviluppo, A/B
  appaiato a $3\times10^{5}$ passi salvo diversa indicazione). I delta
  sono rispetto alla baseline appaiata della loro era e \emph{non} sono
  confrontabili tra versioni dell'ambiente. Dettaglio completo
  nell'\cref{app:registry}.}
  \label{tab:candidates}
  \begin{tabular}{@{}p{4.2cm}p{2.6cm}p{2.2cm}p{3.6cm}@{}}
    \toprule
    Candidato & Tipo & Verdetto & Sintesi \\
    \midrule
    Strumentazione diagnostica & misura & adottato & nessun cambiamento di policy \\
    Obs veicolo route-aware & osservazione & adottato & --- \\
    Blocco di reward shaping (anti-stallo, velocità minima) & reward & promosso & nucleo della \cref{sec:rewards} \\
    Terminazione precoce da stuck & dinamica env & rigettato & SR $-2.9$\,pp, s$+$t $+7.1$\,pp \\
    Correzione lunghezza percorsi & bug fix & adottato & $\approx +50$\,pp SR (cambio di distribuzione) \\
    Correzione seed percorsi & bug fix & adottato & ripristina l'appaiamento A/B \\
    Riparazione value loss & ottimizzatore & mantenuto (gate fallito) & var.\ spiegata $0 \to 0.87$; collisione $+3.3$\,pp \\
    Sconto $0.99 \to 0.997$ & ottimizzatore & mantenuto (falsificato) & timeout $\to$ collisione $\approx$ 1:1 \\
    Schedule di entropia (a frazione) & ottimizzatore & promosso & meccanismo stabile, nessuna instabilità \\
    Penalità di collisione $-50 \to -500$ & reward & mantenuto (falsificato) & collisione piatta ($-0.2$\,pp) \\
    Rimozione della guardia di progresso & reward & promosso & SR $+4.7$\,pp, s$+$t $-3.5$\,pp \\
    Osservabilità 47D (O1$+$O2) & osservazione & adottato & baseline sana, var.\ spiegata 0.983 \\
    Bundle percorso/passo pedoni & reward $+$ percorsi & parziale & passo fuori banda; riflesso sui veicoli \\
    Retuning banda di passo & reward & rigettato & overshoot peggiore; declino in Q4 \\
    Tolleranza gate di hazard $0.75 \to 0.85$ & reward & promosso & 5/5: SR $+2.3$\,pp, s$+$t $-3.0$\,pp, coll./off.\ migliorati \\
    \texttt{route\_short} come successo & misura & revertito & violava il criterio di successo \\
    Tetto attraversamenti $8 \to 3$ & percorsi & promosso & SR $+9.2$\,pp, coll.\ $-2.3$\,pp, off.\ $-2.2$\,pp \\
    \bottomrule
  \end{tabular}
\end{table}

% ------------------------------------------------------------
\section{Design della campagna finale}
\label{sec:campaign}

L'esperimento finale è un fattoriale $2 \times 2$:
\{encoder del critic MLP, GNN\} $\times$ \{curriculum, mixed-batch\} ---
quattro run da $3\times10^{6}$ passi ciascuna, tutte sullo stesso trunk
congelato, con un unico seed di esperimento (999) appaiato tra le
quattro celle, con la difficoltà definita dalla lunghezza del percorso
(\cref{sec:routes}). Le configurazioni delle run sono byte-identiche tra
le celle salvo i due flag sperimentali (regime ed encoder del critic),
proprietà verificata con un diff diretto delle configurazioni. Il design
era originariamente dimensionato a tre seed appaiati per cella; l'asse
di replicazione dei seed è stato eliminato per ragioni di budget di
calcolo (ogni run occupa l'unica macchina disponibile per circa due
giorni, \cref{app:repro}). Questa riduzione è dichiarata qui anziché
nascosta: con un seed per cella, ogni contrasto del \cref{ch:results} è
un \emph{caso di studio appaiato}, non una stima di una distribuzione.

\paragraph{Contrasti.}
Il contrasto primario è curriculum contro mixed-batch \emph{entro}
architettura. Fondamentale: è osservato \emph{due volte} --- una per
ciascuna architettura del critic, addestrata indipendentemente ---
quindi la replicazione tra architetture sostituisce in parte la
replicazione dei seed sull'asse primario: un effetto di regime che
compare con lo stesso segno e struttura comparabile in entrambe le
coppie è molto più difficile da attribuire al rumore da run a run di
quanto lo sarebbe un contrasto singolo. Il contrasto secondario è critic
MLP contro GNN \emph{entro} regime; la lettura dell'interazione
(ipotesi H3, \cref{sec:rq}) è solo direzionale a questa dimensione
campionaria. Nessuna inferenza statistica formale è rivendicata da
nessuna parte: gli effetti sono letti contro la banda delle repliche a
seed identico della \cref{sec:iterations} ($\sim$10\,pp di SR veicoli),
e corroborati dalla coerenza tra livelli di difficoltà, classi di
agenti, scenari di valutazione e strati di misura anziché da test.

\paragraph{Misura.}
Ogni run è misurata ai tre livelli della \cref{sec:eval-pipeline}:
metriche cumulative di addestramento, ricalcolo statico dai log e
valutazione finale del checkpoint su 400 episodi (tre scenari Town03
più lo scenario test Town05 per la generalizzazione di mappa), tutto
riportato per classe.

\paragraph{Scope dichiarato.}
La campagna \emph{non} mette alla prova: famiglie di difficoltà basate
sulla densità di traffico o su assi misti (definite in configurazione ma
inutilizzate); l'encoder del critic ad attenzione e la normalizzazione
PopArt (disponibili, disabilitati); architetture degli actor (sempre
MLP); ablazioni del numero di agenti (sempre $3+3$); o coppie di mappe
oltre Town03 $\to$ Town05.

La matrice delle run realizzata, con verifiche di integrità e stato di
completamento, è riportata nella \cref{tab:runmatrix}.
